% !TEX program = xelatex
% ============================================================
%  Arabic Project Charter Template — Nibras Center
%  مشروع بحثي: ميثاق المشروع
%  Compile with: xelatex main.tex (twice)
% ============================================================

\documentclass[11pt,a4paper]{article}

\usepackage[a4paper,margin=2.5cm,top=2.5cm,bottom=2.5cm,headheight=15pt]{geometry}
\usepackage{amsmath}
\usepackage{graphicx}
\usepackage{array}
\usepackage{booktabs}
\usepackage{tabularx}
\usepackage{multirow}
\usepackage{enumitem}
\usepackage{float}
\usepackage{titlesec}
\usepackage{fancyhdr}
\usepackage{setspace}
\usepackage[table]{xcolor}

\usepackage{bidi}
\usepackage[no-math]{fontspec}
\usepackage[quiet]{polyglossia}

\setmainfont[Scale=1]{NotoKufiArabic-Regular.ttf}
\newfontfamily\arabicfont[Script=Arabic,Scale=0.95]{NotoKufiArabic-Regular.ttf}
\newfontfamily\englishfont[Scale=0.95]{TeX Gyre Termes}
\setmonofont{NotoKufiArabic-Regular.ttf}

\setdefaultlanguage[calendar=gregorian,hijricorrection=1]{arabic}
\setotherlanguage[variant=british]{english}

\usepackage[breaklinks,hidelinks]{hyperref}
\hypersetup{pdftitle={ميثاق المشروع البحثي},pdfauthor={الباحث الرئيسي}}

\addto\captionsarabic{%
  \renewcommand{\contentsname}{الفهرس}%
  \renewcommand{\figurename}{الشكل}%
  \renewcommand{\tablename}{الجدول}%
}

\titleformat{\section}{\Large\bfseries}{\thesection}{0.7em}{}
\titleformat{\subsection}{\large\bfseries}{\thesubsection}{0.6em}{}
\titlespacing*{\section}{0pt}{1.4ex plus 0.4ex}{0.9ex plus 0.3ex}

\pagestyle{fancy}
\fancyhf{}
\fancyhead[R]{\small\itshape ميثاق المشروع}
\fancyhead[L]{\small اسم المشروع}
\fancyfoot[C]{\small\thepage}
\renewcommand{\headrulewidth}{0.3pt}

\setlist[itemize]{topsep=0.3em, itemsep=0.15em, parsep=0pt}
\setlist[enumerate]{topsep=0.3em, itemsep=0.15em, parsep=0pt}

\newcolumntype{L}[1]{>{\raggedright\arraybackslash}p{#1}}
\newcolumntype{C}[1]{>{\centering\arraybackslash}p{#1}}

\definecolor{headerColor}{HTML}{1A3C6B}
\definecolor{lightRow}{HTML}{F0F4F8}

\setstretch{1.4}

\begin{document}

% ============================================================
% TITLE BLOCK
% ============================================================
\begin{center}
{\LARGE\bfseries ميثاق المشروع البحثي}\\[0.5em]
{\Large عنوان المشروع البحثي}\\[1em]
\end{center}

% ---- Project info table ----
\begin{table}[H]
\centering
\renewcommand{\arraystretch}{1.4}
\begin{tabularx}{\textwidth}{L{4cm} X L{4cm} X}
\toprule
\textbf{اسم المشروع:} & \multicolumn{3}{l}{---} \\
\textbf{رقم المشروع:} & --- & \textbf{تاريخ الإصدار:} & --- \\
\textbf{الباحث الرئيسي:} & --- & \textbf{الجهة الممولة:} & --- \\
\textbf{الجامعة:} & --- & \textbf{القسم:} & --- \\
\textbf{تاريخ البدء:} & --- & \textbf{تاريخ الانتهاء:} & --- \\
\textbf{الميزانية:} & --- & \textbf{المدة:} & --- \\
\bottomrule
\end{tabularx}
\end{table}

\vspace{1em}

% ============================================================
% 1. PROJECT OVERVIEW
% ============================================================
\section{نظرة عامة على المشروع}
\label{sec:overview}

% TODO: اكتب نظرة عامة مختصرة على المشروع (فقرة واحدة).
هذا نص تجريبي للنظرة العامة. يقدم وصفاً مختصراً للمشروع البحثي وأهميته.

% ============================================================
% 2. OBJECTIVES
% ============================================================
\section{أهداف المشروع}
\label{sec:objectives}

% TODO: اكتب أهداف المشروع.
\begin{enumerate}
    \item الهدف الأول: وصف الهدف.
    \item الهدف الثاني: وصف الهدف.
    \item الهدف الثالث: وصف الهدف.
\end{enumerate}

% ============================================================
% 3. SCOPE
% ============================================================
\section{نطاق المشروع}
\label{sec:scope}

\subsection{المشمولات}
% TODO: اذكر ما يشمله المشروع.
\begin{itemize}
    \item البند الأول المشمول في المشروع.
    \item البند الثاني المشمول في المشروع.
\end{itemize}

\subsection{غير المشمولات}
% TODO: اذكر ما لا يشمله المشروع.
\begin{itemize}
    \item البند الأول غير المشمول.
    \item البند الثاني غير المشمول.
\end{itemize}

% ============================================================
% 4. STAKEHOLDERS
% ============================================================
\section{أصحاب المصلحة}
\label{sec:stakeholders}

\begin{table}[H]
\centering
\caption{أصحاب المصلحة الرئيسيون}
\renewcommand{\arraystretch}{1.3}
\begin{tabularx}{\textwidth}{L{4cm} L{4cm} X}
\toprule
\textbf{الاسم} & \textbf{الدور} & \textbf{المسؤوليات} \\
\midrule
--- & الباحث الرئيسي & الإشراف العام وإدارة المشروع \\
--- & باحث مشارك & تنفيذ التجارب وتحليل البيانات \\
--- & الجهة الممولة & تمويل المشروع ومتابعة التقدم \\
--- & المستفيدون & الاستفادة من نتائج المشروع \\
\bottomrule
\end{tabularx}
\end{table}

% ============================================================
% 5. DELIVERABLES
% ============================================================
\section{المخرجات الرئيسية}
\label{sec:deliverables}

\begin{table}[H]
\centering
\caption{المخرجات الرئيسية للمشروع}
\renewcommand{\arraystretch}{1.3}
\begin{tabularx}{\textwidth}{C{1cm} L{5cm} L{3cm} X}
\toprule
\textbf{رقم} & \textbf{المخرج} & \textbf{تاريخ التسليم} & \textbf{معايير القبول} \\
\midrule
1 & تقرير المراجعة والأدبيات & --- & مراجعة شاملة ومحكمة \\
2 & تصميم التجربة & --- & موافقة اللجنة العلمية \\
3 & جمع البيانات & --- & اكتمال البيانات لجميع العينات \\
4 & التحليل والنتائج & --- & دقة إحصائية مقبولة \\
5 & الورقة البحثية & --- & قبولها للنشر في مجلة محكمة \\
6 & التقرير النهائي & --- & موافقة الجهة الممولة \\
\bottomrule
\end{tabularx}
\end{table}

% ============================================================
% 6. MILESTONES
% ============================================================
\section{المعالم الرئيسية (Milestones)}
\label{sec:milestones}

\begin{table}[H]
\centering
\caption{المعالم الرئيسية للمشروع}
\renewcommand{\arraystretch}{1.3}
\begin{tabularx}{\textwidth}{C{1cm} L{6cm} L{4cm} X}
\toprule
\textbf{رقم} & \textbf{المعلم} & \textbf{التاريخ المتوقع} & \textbf{الحالة} \\
\midrule
M1 & انتهاء المراجعة والأدبيات & --- & لم يبدأ \\
M2 & انتهاء تصميم التجربة & --- & لم يبدأ \\
M3 & انتهاء جمع البيانات & --- & لم يبدأ \\
M4 & انتهاء التحليل & --- & لم يبدأ \\
M5 & تسليم الورقة البحثية & --- & لم يبدأ \\
M6 & التقرير النهائي وإغلاق المشروع & --- & لم يبدأ \\
\bottomrule
\end{tabularx}
\end{table}

% ============================================================
% 7. BUDGET SUMMARY
% ============================================================
\section{ملخص الميزانية}
\label{sec:budget}

\begin{table}[H]
\centering
\caption{ملخص الميزانية}
\begin{tabular}{|l|r|}
\hline
\textbf{البند} & \textbf{المبلغ (\LR{USD})} \\
\hline
الرواتب والأجور & 0 \\
المعدات & 0 \\
المواد الاستهلاكية & 0 \\
السفر والمؤتمرات & 0 \\
النشر & 0 \\
احتياطي طوارئ (10\%) & 0 \\
\hline
\textbf{المجموع} & \textbf{0} \\
\hline
\end{tabular}
\end{table}

% ============================================================
% 8. RISKS SUMMARY
% ============================================================
\section{ملخص المخاطر الرئيسية}
\label{sec:risks}

\begin{table}[H]
\centering
\caption{أبرز المخاطر}
\renewcommand{\arraystretch}{1.3}
\begin{tabularx}{\textwidth}{C{1cm} L{4cm} C{2cm} C{2cm} X}
\toprule
\textbf{رقم} & \textbf{المخاطرة} & \textbf{الاحتمال} & \textbf{التأثير} & \textbf{استراتيجية المعالجة} \\
\midrule
1 & تأخر في جمع البيانات & متوسط & عالٍ & خطة بديلة للمصادر \\
2 & فشل الأجهزة & منخفض & عالٍ & صيانة دورية وقطع غيار \\
3 & رفض الورقة للنشر & متوسط & متوسط & مراجعة قبل التقديم \\
\bottomrule
\end{tabularx}
\end{table}

% ============================================================
% 9. APPROVALS
% ============================================================
\section{الموافقات والاعتمادات}
\label{sec:approvals}

\vspace{2em}
\begin{table}[H]
\centering
\renewcommand{\arraystretch}{2}
\begin{tabularx}{\textwidth}{L{4cm} L{5cm} X}
\toprule
\textbf{الدور} & \textbf{الاسم} & \textbf{التوقيع والتاريخ} \\
\midrule
الباحث الرئيسي & --- & \rule{6cm}{0.4pt} \\
مدير المشروع & --- & \rule{6cm}{0.4pt} \\
ممثل الجهة الممولة & --- & \rule{6cm}{0.4pt} \\
\bottomrule
\end{tabularx}
\end{table}

\end{document}
